% Options for packages loaded elsewhere
\PassOptionsToPackage{unicode}{hyperref}
\PassOptionsToPackage{hyphens}{url}
\PassOptionsToPackage{dvipsnames,svgnames,x11names}{xcolor}
\documentclass[
  11pt,
]{article}
\usepackage{xcolor}
\usepackage[margin=0.5in]{geometry}
\usepackage{amsmath,amssymb}
\setcounter{secnumdepth}{5}
\usepackage{iftex}
\ifPDFTeX
  \usepackage[T1]{fontenc}
  \usepackage[utf8]{inputenc}
  \usepackage{textcomp} % provide euro and other symbols
\else % if luatex or xetex
  \usepackage{unicode-math} % this also loads fontspec
  \defaultfontfeatures{Scale=MatchLowercase}
  \defaultfontfeatures[\rmfamily]{Ligatures=TeX,Scale=1}
\fi
\usepackage{lmodern}
\ifPDFTeX\else
  % xetex/luatex font selection
  \setmainfont[]{Helvetica}
  \setmonofont[]{Menlo}
\fi
% Use upquote if available, for straight quotes in verbatim environments
\IfFileExists{upquote.sty}{\usepackage{upquote}}{}
\IfFileExists{microtype.sty}{% use microtype if available
  \usepackage[]{microtype}
  \UseMicrotypeSet[protrusion]{basicmath} % disable protrusion for tt fonts
}{}
\makeatletter
\@ifundefined{KOMAClassName}{% if non-KOMA class
  \IfFileExists{parskip.sty}{%
    \usepackage{parskip}
  }{% else
    \setlength{\parindent}{0pt}
    \setlength{\parskip}{6pt plus 2pt minus 1pt}}
}{% if KOMA class
  \KOMAoptions{parskip=half}}
\makeatother
\usepackage{color}
\usepackage{fancyvrb}
\newcommand{\VerbBar}{|}
\newcommand{\VERB}{\Verb[commandchars=\\\{\}]}
\DefineVerbatimEnvironment{Highlighting}{Verbatim}{commandchars=\\\{\}}
% Add ',fontsize=\small' for more characters per line
\usepackage{framed}
\definecolor{shadecolor}{RGB}{248,248,248}
\newenvironment{Shaded}{\begin{snugshade}}{\end{snugshade}}
\newcommand{\AlertTok}[1]{\textcolor[rgb]{0.94,0.16,0.16}{#1}}
\newcommand{\AnnotationTok}[1]{\textcolor[rgb]{0.56,0.35,0.01}{\textbf{\textit{#1}}}}
\newcommand{\AttributeTok}[1]{\textcolor[rgb]{0.13,0.29,0.53}{#1}}
\newcommand{\BaseNTok}[1]{\textcolor[rgb]{0.00,0.00,0.81}{#1}}
\newcommand{\BuiltInTok}[1]{#1}
\newcommand{\CharTok}[1]{\textcolor[rgb]{0.31,0.60,0.02}{#1}}
\newcommand{\CommentTok}[1]{\textcolor[rgb]{0.56,0.35,0.01}{\textit{#1}}}
\newcommand{\CommentVarTok}[1]{\textcolor[rgb]{0.56,0.35,0.01}{\textbf{\textit{#1}}}}
\newcommand{\ConstantTok}[1]{\textcolor[rgb]{0.56,0.35,0.01}{#1}}
\newcommand{\ControlFlowTok}[1]{\textcolor[rgb]{0.13,0.29,0.53}{\textbf{#1}}}
\newcommand{\DataTypeTok}[1]{\textcolor[rgb]{0.13,0.29,0.53}{#1}}
\newcommand{\DecValTok}[1]{\textcolor[rgb]{0.00,0.00,0.81}{#1}}
\newcommand{\DocumentationTok}[1]{\textcolor[rgb]{0.56,0.35,0.01}{\textbf{\textit{#1}}}}
\newcommand{\ErrorTok}[1]{\textcolor[rgb]{0.64,0.00,0.00}{\textbf{#1}}}
\newcommand{\ExtensionTok}[1]{#1}
\newcommand{\FloatTok}[1]{\textcolor[rgb]{0.00,0.00,0.81}{#1}}
\newcommand{\FunctionTok}[1]{\textcolor[rgb]{0.13,0.29,0.53}{\textbf{#1}}}
\newcommand{\ImportTok}[1]{#1}
\newcommand{\InformationTok}[1]{\textcolor[rgb]{0.56,0.35,0.01}{\textbf{\textit{#1}}}}
\newcommand{\KeywordTok}[1]{\textcolor[rgb]{0.13,0.29,0.53}{\textbf{#1}}}
\newcommand{\NormalTok}[1]{#1}
\newcommand{\OperatorTok}[1]{\textcolor[rgb]{0.81,0.36,0.00}{\textbf{#1}}}
\newcommand{\OtherTok}[1]{\textcolor[rgb]{0.56,0.35,0.01}{#1}}
\newcommand{\PreprocessorTok}[1]{\textcolor[rgb]{0.56,0.35,0.01}{\textit{#1}}}
\newcommand{\RegionMarkerTok}[1]{#1}
\newcommand{\SpecialCharTok}[1]{\textcolor[rgb]{0.81,0.36,0.00}{\textbf{#1}}}
\newcommand{\SpecialStringTok}[1]{\textcolor[rgb]{0.31,0.60,0.02}{#1}}
\newcommand{\StringTok}[1]{\textcolor[rgb]{0.31,0.60,0.02}{#1}}
\newcommand{\VariableTok}[1]{\textcolor[rgb]{0.00,0.00,0.00}{#1}}
\newcommand{\VerbatimStringTok}[1]{\textcolor[rgb]{0.31,0.60,0.02}{#1}}
\newcommand{\WarningTok}[1]{\textcolor[rgb]{0.56,0.35,0.01}{\textbf{\textit{#1}}}}
\usepackage{longtable,booktabs,array}
\newcounter{none} % for unnumbered tables
\usepackage{calc} % for calculating minipage widths
% Correct order of tables after \paragraph or \subparagraph
\usepackage{etoolbox}
\makeatletter
\patchcmd\longtable{\par}{\if@noskipsec\mbox{}\fi\par}{}{}
\makeatother
% Allow footnotes in longtable head/foot
\IfFileExists{footnotehyper.sty}{\usepackage{footnotehyper}}{\usepackage{footnote}}
\makesavenoteenv{longtable}
\usepackage{graphicx}
\makeatletter
\newsavebox\pandoc@box
\newcommand*\pandocbounded[1]{% scales image to fit in text height/width
  \sbox\pandoc@box{#1}%
  \Gscale@div\@tempa{\textheight}{\dimexpr\ht\pandoc@box+\dp\pandoc@box\relax}%
  \Gscale@div\@tempb{\linewidth}{\wd\pandoc@box}%
  \ifdim\@tempb\p@<\@tempa\p@\let\@tempa\@tempb\fi% select the smaller of both
  \ifdim\@tempa\p@<\p@\scalebox{\@tempa}{\usebox\pandoc@box}%
  \else\usebox{\pandoc@box}%
  \fi%
}
% Set default figure placement to htbp
\def\fps@figure{htbp}
\makeatother
\setlength{\emergencystretch}{3em} % prevent overfull lines
\providecommand{\tightlist}{%
  \setlength{\itemsep}{0pt}\setlength{\parskip}{0pt}}
\usepackage{xcolor}
\usepackage{fvextra}
\DefineVerbatimEnvironment{Highlighting}{Verbatim}{breaklines,breakanywhere,commandchars=\\\{\}}
\usepackage{graphicx}
\usepackage{float}
\usepackage{microtype}
\usepackage{setspace}
\usepackage{bookmark}
\IfFileExists{xurl.sty}{\usepackage{xurl}}{} % add URL line breaks if available
\urlstyle{same}
\hypersetup{
  pdftitle={624 Trees: Learning from the Central Park Squirrel Census},
  pdfauthor={Rebecca Bronstein},
  colorlinks=true,
  linkcolor={blue},
  filecolor={Maroon},
  citecolor={Blue},
  urlcolor={blue},
  pdfcreator={LaTeX via pandoc}}

\title{624 Trees: Learning from the Central Park Squirrel Census}
\author{Rebecca Bronstein}
\date{2026-03-30}

\begin{document}
\maketitle

{
\setcounter{tocdepth}{2}
\tableofcontents
}
\subsection{Why Use Tree Methods on Squirrel
Data?}\label{why-use-tree-methods-on-squirrel-data}

Tree models are helpful here because the squirrel census is full of
mixed data types: categorical fur colors, yes/no behaviors, vague human
interaction signals, and spatial labels such as hectares. A single tree
can turn that messy field data into readable splits, while ensembles let
us check whether the same patterns hold up when we average across many
trees.

Instead of treating this as a generic modeling exercise, I use tree
methods as a field guide:

\begin{enumerate}
\def\labelenumi{\arabic{enumi}.}
\tightlist
\item
  A classification tree asks where squirrels are likely to be observed:
  on the ground or above it.
\item
  A regression tree asks how high squirrels perch once they are already
  above ground.
\item
  A rule-based tree turns feeding behavior into short, interpretable
  rules.
\item
  A bagged tree checks whether an ensemble improves the stability of the
  location prediction.
\end{enumerate}

\subsection{Data Snapshot}\label{data-snapshot}

\begin{Shaded}
\begin{Highlighting}[]
\FunctionTok{kable}\NormalTok{(overview, }\AttributeTok{caption =} \StringTok{"Quick overview of the squirrel census data used in this report."}\NormalTok{)}
\end{Highlighting}
\end{Shaded}

\begin{longtable}[]{@{}lr@{}}
\caption{Quick overview of the squirrel census data used in this
report.}\tabularnewline
\toprule\noalign{}
Measure & Value \\
\midrule\noalign{}
\endfirsthead
\toprule\noalign{}
Measure & Value \\
\midrule\noalign{}
\endhead
\bottomrule\noalign{}
\endlastfoot
Total observations & 3023 \\
Distinct hectares & 339 \\
Ground-plane sightings & 2116 \\
Above-ground sightings & 843 \\
Sightings with recorded height & 793 \\
\end{longtable}

\subsection{Classification Tree: What Predicts Ground vs.~Above-Ground
Sightings?}\label{classification-tree-what-predicts-ground-vs.-above-ground-sightings}

This tree treats the park like a behavioral landscape. The goal is to
see whether movement, feeding, vocalization, tail signals, and basic
location cues help separate squirrels seen on the ground from squirrels
spotted above ground.

\begin{Shaded}
\begin{Highlighting}[]
\NormalTok{location\_tree }\OtherTok{\textless{}{-}} \FunctionTok{rpart}\NormalTok{(Location.Class }\SpecialCharTok{\textasciitilde{}}\NormalTok{ ., }\AttributeTok{data =}\NormalTok{ train\_location, }\AttributeTok{method =} \StringTok{"class"}\NormalTok{)}
\FunctionTok{plot}\NormalTok{(location\_tree, }\AttributeTok{uniform =} \ConstantTok{TRUE}\NormalTok{, }\AttributeTok{margin =} \FloatTok{0.08}\NormalTok{)}
\FunctionTok{text}\NormalTok{(location\_tree, }\AttributeTok{use.n =} \ConstantTok{TRUE}\NormalTok{, }\AttributeTok{cex =} \FloatTok{0.7}\NormalTok{)}
\end{Highlighting}
\end{Shaded}

\pandocbounded{\includegraphics[keepaspectratio]{trees_files/figure-latex/classification_tree-1.pdf}}

\begin{Shaded}
\begin{Highlighting}[]
\NormalTok{location\_pred }\OtherTok{\textless{}{-}} \FunctionTok{predict}\NormalTok{(location\_tree, test\_location, }\AttributeTok{type =} \StringTok{"class"}\NormalTok{)}
\NormalTok{location\_metrics }\OtherTok{\textless{}{-}} \FunctionTok{classification\_metrics}\NormalTok{(test\_location}\SpecialCharTok{$}\NormalTok{Location.Class, location\_pred)}
\NormalTok{location\_importance }\OtherTok{\textless{}{-}} \FunctionTok{top\_importance}\NormalTok{(location\_tree)}

\FunctionTok{kable}\NormalTok{(location\_metrics, }\AttributeTok{caption =} \StringTok{"Classification tree performance for predicting squirrel location."}\NormalTok{)}
\end{Highlighting}
\end{Shaded}

\begin{longtable}[]{@{}rr@{}}
\caption{Classification tree performance for predicting squirrel
location.}\tabularnewline
\toprule\noalign{}
Accuracy & Kappa \\
\midrule\noalign{}
\endfirsthead
\toprule\noalign{}
Accuracy & Kappa \\
\midrule\noalign{}
\endhead
\bottomrule\noalign{}
\endlastfoot
0.749 & 0.271 \\
\end{longtable}

\begin{Shaded}
\begin{Highlighting}[]
\FunctionTok{kable}\NormalTok{(location\_importance, }\AttributeTok{caption =} \StringTok{"Most influential variables in the location tree."}\NormalTok{)}
\end{Highlighting}
\end{Shaded}

\begin{longtable}[]{@{}lr@{}}
\caption{Most influential variables in the location
tree.}\tabularnewline
\toprule\noalign{}
Variable & Importance \\
\midrule\noalign{}
\endfirsthead
\toprule\noalign{}
Variable & Importance \\
\midrule\noalign{}
\endhead
\bottomrule\noalign{}
\endlastfoot
Feeding & 121.25 \\
Active & 56.93 \\
Human.Response & 12.43 \\
Shift & 9.36 \\
Vocal & 6.21 \\
Hectare.Number & 5.25 \\
\end{longtable}

The classification tree helps answer a practical question for field
observation: what makes a squirrel more likely to be seen above ground?
In this data, active movement and feeding-related behavior matter much
more than fur color, which suggests behavior is more informative than
appearance for understanding where squirrels position themselves.

\subsection{Regression Tree: How High Do Above-Ground Squirrels
Perch?}\label{regression-tree-how-high-do-above-ground-squirrels-perch}

Once a squirrel is above ground, the next question is not just
\emph{whether} it climbed, but \emph{how high}. This regression tree
uses the same mix of spatial and behavioral cues to estimate recorded
height.

\begin{Shaded}
\begin{Highlighting}[]
\NormalTok{height\_tree }\OtherTok{\textless{}{-}} \FunctionTok{rpart}\NormalTok{(Above.Ground.Height }\SpecialCharTok{\textasciitilde{}}\NormalTok{ ., }\AttributeTok{data =}\NormalTok{ train\_height, }\AttributeTok{method =} \StringTok{"anova"}\NormalTok{)}
\FunctionTok{plot}\NormalTok{(height\_tree, }\AttributeTok{uniform =} \ConstantTok{TRUE}\NormalTok{, }\AttributeTok{margin =} \FloatTok{0.08}\NormalTok{)}
\FunctionTok{text}\NormalTok{(height\_tree, }\AttributeTok{use.n =} \ConstantTok{TRUE}\NormalTok{, }\AttributeTok{cex =} \FloatTok{0.7}\NormalTok{)}
\end{Highlighting}
\end{Shaded}

\pandocbounded{\includegraphics[keepaspectratio]{trees_files/figure-latex/regression_tree-1.pdf}}

\begin{Shaded}
\begin{Highlighting}[]
\NormalTok{height\_pred }\OtherTok{\textless{}{-}} \FunctionTok{predict}\NormalTok{(height\_tree, test\_height)}
\NormalTok{height\_metrics }\OtherTok{\textless{}{-}} \FunctionTok{regression\_metrics}\NormalTok{(test\_height}\SpecialCharTok{$}\NormalTok{Above.Ground.Height, height\_pred)}
\NormalTok{height\_importance }\OtherTok{\textless{}{-}} \FunctionTok{top\_importance}\NormalTok{(height\_tree)}

\FunctionTok{kable}\NormalTok{(height\_metrics, }\AttributeTok{caption =} \StringTok{"Regression tree performance for above{-}ground height."}\NormalTok{)}
\end{Highlighting}
\end{Shaded}

\begin{longtable}[]{@{}rr@{}}
\caption{Regression tree performance for above-ground
height.}\tabularnewline
\toprule\noalign{}
RMSE & MAE \\
\midrule\noalign{}
\endfirsthead
\toprule\noalign{}
RMSE & MAE \\
\midrule\noalign{}
\endhead
\bottomrule\noalign{}
\endlastfoot
12.42 & 9.11 \\
\end{longtable}

\begin{Shaded}
\begin{Highlighting}[]
\FunctionTok{kable}\NormalTok{(height\_importance, }\AttributeTok{caption =} \StringTok{"Most influential variables in the height tree."}\NormalTok{)}
\end{Highlighting}
\end{Shaded}

\begin{longtable}[]{@{}lr@{}}
\caption{Most influential variables in the height tree.}\tabularnewline
\toprule\noalign{}
Variable & Importance \\
\midrule\noalign{}
\endfirsthead
\toprule\noalign{}
Variable & Importance \\
\midrule\noalign{}
\endhead
\bottomrule\noalign{}
\endlastfoot
Hectare.Number & 10877.62 \\
Human.Response & 8211.92 \\
Hectare.Letter & 6634.14 \\
Shift & 1497.91 \\
Fur.Color & 1370.95 \\
Tail.Signal & 331.20 \\
\end{longtable}

This model treats height as a kind of vertical habitat choice. The most
important predictors are park position and human-response context rather
than simple color categories, which hints that where a squirrel is seen
in the park may be more useful than how it looks when we want to explain
height.

\subsection{Rule-Based Tree: What Simple Rules Describe Feeding
Behavior?}\label{rule-based-tree-what-simple-rules-describe-feeding-behavior}

Rule models are useful when we want a short, readable description
instead of a large branching diagram. Here, I predict whether a squirrel
is feeding or not feeding from its activity level, location, and other
observation cues.

\begin{Shaded}
\begin{Highlighting}[]
\NormalTok{feeding\_rules }\OtherTok{\textless{}{-}} \FunctionTok{C5.0}\NormalTok{(Feeding }\SpecialCharTok{\textasciitilde{}}\NormalTok{ ., }\AttributeTok{data =}\NormalTok{ train\_feeding, }\AttributeTok{rules =} \ConstantTok{TRUE}\NormalTok{)}
\NormalTok{feeding\_pred }\OtherTok{\textless{}{-}} \FunctionTok{predict}\NormalTok{(feeding\_rules, test\_feeding)}
\NormalTok{feeding\_metrics }\OtherTok{\textless{}{-}} \FunctionTok{classification\_metrics}\NormalTok{(test\_feeding}\SpecialCharTok{$}\NormalTok{Feeding, feeding\_pred)}
\NormalTok{feeding\_usage }\OtherTok{\textless{}{-}} \FunctionTok{as.data.frame}\NormalTok{(}\FunctionTok{C5imp}\NormalTok{(feeding\_rules, }\AttributeTok{metric =} \StringTok{"usage"}\NormalTok{)) }\SpecialCharTok{\%\textgreater{}\%}
  \FunctionTok{rownames\_to\_column}\NormalTok{(}\AttributeTok{var =} \StringTok{"Variable"}\NormalTok{) }\SpecialCharTok{\%\textgreater{}\%}
  \FunctionTok{rename}\NormalTok{(}\AttributeTok{Usage =}\NormalTok{ Overall) }\SpecialCharTok{\%\textgreater{}\%}
  \FunctionTok{arrange}\NormalTok{(}\FunctionTok{desc}\NormalTok{(Usage)) }\SpecialCharTok{\%\textgreater{}\%}
  \FunctionTok{mutate}\NormalTok{(}\AttributeTok{Usage =} \FunctionTok{round}\NormalTok{(Usage, }\DecValTok{2}\NormalTok{))}

\FunctionTok{kable}\NormalTok{(feeding\_metrics, }\AttributeTok{caption =} \StringTok{"Rule{-}based tree performance for feeding behavior."}\NormalTok{)}
\end{Highlighting}
\end{Shaded}

\begin{longtable}[]{@{}rr@{}}
\caption{Rule-based tree performance for feeding
behavior.}\tabularnewline
\toprule\noalign{}
Accuracy & Kappa \\
\midrule\noalign{}
\endfirsthead
\toprule\noalign{}
Accuracy & Kappa \\
\midrule\noalign{}
\endhead
\bottomrule\noalign{}
\endlastfoot
0.778 & 0.55 \\
\end{longtable}

\begin{Shaded}
\begin{Highlighting}[]
\FunctionTok{kable}\NormalTok{(}\FunctionTok{head}\NormalTok{(feeding\_usage, }\DecValTok{6}\NormalTok{), }\AttributeTok{caption =} \StringTok{"Most{-}used predictors in the feeding rules."}\NormalTok{)}
\end{Highlighting}
\end{Shaded}

\begin{longtable}[]{@{}lr@{}}
\caption{Most-used predictors in the feeding rules.}\tabularnewline
\toprule\noalign{}
Variable & Usage \\
\midrule\noalign{}
\endfirsthead
\toprule\noalign{}
Variable & Usage \\
\midrule\noalign{}
\endhead
\bottomrule\noalign{}
\endlastfoot
Active & 100 \\
Shift & 0 \\
Age & 0 \\
Fur.Color & 0 \\
Hectare.Number & 0 \\
Hectare.Letter & 0 \\
\end{longtable}

\begin{Shaded}
\begin{Highlighting}[]
\FunctionTok{summary}\NormalTok{(feeding\_rules)}
\end{Highlighting}
\end{Shaded}

\begin{verbatim}
## 
## Call:
## C5.0.formula(formula = Feeding ~ ., data = train_feeding, rules = TRUE)
## 
## 
## C5.0 [Release 2.07 GPL Edition]      Wed Apr  1 21:30:42 2026
## -------------------------------
## 
## Class specified by attribute `outcome'
## 
## Read 1976 cases (11 attributes) from undefined.data
## 
## Rules:
## 
## Rule 1: (1019/102, lift 1.5)
##  Active = Still
##  ->  class Feeding  [0.899]
## 
## Rule 2: (957/298, lift 1.8)
##  Active = Active
##  ->  class Not feeding  [0.688]
## 
## Default class: Feeding
## 
## 
## Evaluation on training data (1976 cases):
## 
##          Rules     
##    ----------------
##      No      Errors
## 
##       2  400(20.2%)   <<
## 
## 
##     (a)   (b)    <-classified as
##    ----  ----
##     917   298    (a): class Feeding
##     102   659    (b): class Not feeding
## 
## 
##  Attribute usage:
## 
##  100.00% Active
## 
## 
## Time: 0.0 secs
\end{verbatim}

The rules are surprisingly compact. In this dataset, activity level does
a lot of the work: squirrels that are still are much more often recorded
as feeding, while squirrels that are actively running, climbing, or
chasing are more often recorded as not feeding. That is exactly the kind
of ecological shorthand tree methods are good at surfacing.

\subsection{Bagged Trees: Does an Ensemble Improve the Location
Story?}\label{bagged-trees-does-an-ensemble-improve-the-location-story}

A single tree is easy to interpret, but it can be sensitive to how the
sample is split. Bagging averages over many trees, which usually gives a
more stable answer. I apply bagging to the same ground-vs-above-ground
problem used in the first section.

\begin{Shaded}
\begin{Highlighting}[]
\NormalTok{bag\_model }\OtherTok{\textless{}{-}} \FunctionTok{bagging}\NormalTok{(Location.Class }\SpecialCharTok{\textasciitilde{}}\NormalTok{ ., }\AttributeTok{data =}\NormalTok{ train\_location, }\AttributeTok{nbagg =} \DecValTok{25}\NormalTok{)}
\NormalTok{bag\_pred }\OtherTok{\textless{}{-}} \FunctionTok{predict}\NormalTok{(bag\_model, test\_location)}

\ControlFlowTok{if}\NormalTok{ (}\FunctionTok{is.list}\NormalTok{(bag\_pred) }\SpecialCharTok{\&\&} \StringTok{"class"} \SpecialCharTok{\%in\%} \FunctionTok{names}\NormalTok{(bag\_pred)) \{}
\NormalTok{  bag\_pred }\OtherTok{\textless{}{-}}\NormalTok{ bag\_pred}\SpecialCharTok{$}\NormalTok{class}
\NormalTok{\}}

\NormalTok{bag\_pred }\OtherTok{\textless{}{-}} \FunctionTok{factor}\NormalTok{(bag\_pred, }\AttributeTok{levels =} \FunctionTok{levels}\NormalTok{(test\_location}\SpecialCharTok{$}\NormalTok{Location.Class))}

\NormalTok{bag\_metrics }\OtherTok{\textless{}{-}} \FunctionTok{classification\_metrics}\NormalTok{(test\_location}\SpecialCharTok{$}\NormalTok{Location.Class, bag\_pred) }\SpecialCharTok{\%\textgreater{}\%}
  \FunctionTok{mutate}\NormalTok{(}\AttributeTok{Model =} \StringTok{"Bagged trees"}\NormalTok{)}

\NormalTok{single\_tree\_metrics }\OtherTok{\textless{}{-}}\NormalTok{ location\_metrics }\SpecialCharTok{\%\textgreater{}\%}
  \FunctionTok{mutate}\NormalTok{(}\AttributeTok{Model =} \StringTok{"Single tree"}\NormalTok{)}

\NormalTok{comparison }\OtherTok{\textless{}{-}} \FunctionTok{bind\_rows}\NormalTok{(single\_tree\_metrics, bag\_metrics) }\SpecialCharTok{\%\textgreater{}\%}
  \FunctionTok{select}\NormalTok{(Model, Accuracy, Kappa)}

\FunctionTok{kable}\NormalTok{(comparison, }\AttributeTok{caption =} \StringTok{"Comparing a single location tree with a bagged ensemble."}\NormalTok{)}
\end{Highlighting}
\end{Shaded}

\begin{longtable}[]{@{}lrr@{}}
\caption{Comparing a single location tree with a bagged
ensemble.}\tabularnewline
\toprule\noalign{}
Model & Accuracy & Kappa \\
\midrule\noalign{}
\endfirsthead
\toprule\noalign{}
Model & Accuracy & Kappa \\
\midrule\noalign{}
\endhead
\bottomrule\noalign{}
\endlastfoot
Single tree & 0.749 & 0.271 \\
Bagged trees & 0.739 & 0.319 \\
\end{longtable}

The bagged model is only slightly more accurate here, but that
comparison is still useful. It suggests that the strongest patterns in
this squirrel census are already fairly stable, so the simple tree
captures much of the signal while the ensemble adds a small amount of
extra reliability.

\subsection{Conclusion}\label{conclusion}

Tree methods help turn this squirrel census into more than a table of
observations. They show that:

\begin{enumerate}
\def\labelenumi{\arabic{enumi}.}
\tightlist
\item
  Behavior is more informative than fur color for predicting where
  squirrels are seen.
\item
  Above-ground height looks partly like a spatial decision, not just a
  random measurement.
\item
  Feeding behavior can be summarized with a small set of readable rules.
\item
  Bagging confirms that the main location patterns are real enough to
  survive beyond a single tree.
\end{enumerate}

That is the main value of tree-based models in a dataset like this one:
they do not just predict outcomes, they organize messy field
observations into a form that is easier to explain.

\end{document}
